% SOURCE_AUTHORITY: sga5_fr_workpass.tex, lines 10343--10386 (LF-normalized unit).
% SOURCE_SHA256: 791F4EFFC5E02832D5D77ED03518C8156D6F07E4C8238B03545DB93D883FBB28
\section*{7. Intersecciones completas}

En este apartado se supone dado un cuerpo separablemente cerrado $k$. Si $\ell$ es un número primo $\ne\operatorname{car}(k)$, para todo $k$-esquema propio $X$ y todo entero $i\ge0$, se denota por
\[
b_\ell^i(X)\quad\text{o}\quad b^i(X)
\]
si no hay posibilidad de confusión, el $i$-ésimo número de Betti de $X$ con coeficientes en $\mathbb Z_\ell$. Si $X$ es proyectivo sobre $k$, se llama sección hiperplana de $X$ al lugar de los ceros de una sección transversal a la sección cero de un fibrado vectorial $V(\check{\mathcal L})$, donde $\mathcal L$ es un $\mathcal O_X$-módulo invertible amplio.

\begin{statement}{Teorema 7.1}
(Teorema de Lefschetz.) Sean $X$ un $k$-esquema proyectivo liso, conexo, de dimensión $n$, $Y$ un subesquema de $X$ que es intersección de $d$ secciones hiperplanas de $X$, y $L$ un $A_X$-módulo localmente constante constructible. Entonces:
\begin{enumerate}[label=(\roman*)]
\item El homomorfismo de imagen inversa
\[
H^i(X,L)\to H^i(Y,L|Y)
\]
es un isomorfismo para $i\le n-d-1$, y un monomorfismo para $i=n-d$.
\item Supongamos que $Y$ es liso sobre $k$, de dimensión $n-d$. Entonces el homomorfismo de Gysin
\[
H^i(Y,(L|Y)\otimes\mu_Y^{\otimes -d})\to H^{i+2d}(X,L)
\]
es un isomorfismo para $i\ge n-d+1$ y un epimorfismo para $i=n-d$.
\end{enumerate}
\end{statement}

\emph{Demostración.} Es bien sabido que el complementario de una sección hiperplana es un abierto afín y, por consiguiente, $U=X-Y$ es unión de $d$ abiertos afines. La afirmación (i) se deduce inmediatamente de la sucesión exacta de cohomología
\[
\cdots\to H_!^i(U,L|U)\to H^i(X,L)\to H^i(Y,L|Y)
\xrightarrow{\partial} H_!^{i+1}(U,L|U)\to\cdots
\]
y del lema siguiente.

\begin{statement}{Lema 7.1.1}
Sea $U$ un $k$-esquema de tipo finito de dimensión $n$, unión de $d$ abiertos afines, y $F$ un $A_U$-módulo. Entonces:
\begin{enumerate}[label=\alph*)]
\item si $F$ es constructible,
\[
H^i(U,F)=0\qquad\text{para }i\ge n+d;
\]
\item si $F$ es localmente constante y $U$ es liso sobre $k$,
\[
H^i_c(U,F)=0\qquad\text{para }i\le n-d.
\]
\end{enumerate}
\end{statement}
